%\documentclass[11pt,a4paper,twoside]{article}

\usepackage[utf8]{inputenc}

\usepackage[T1,T2A]{fontenc}

\usepackage[english.ancient,russian.ancient]{babel}

\usepackage{graphicx}

\usepackage{array}

\usepackage[doi]{samgtu-bib}

\usepackage{samgtu}

\usepackage{samgtu-name}

\usepackage{mathtools}

\mathtoolsset{showonlyrefs}

\usepackage{desclist}

\usepackage{euscript}

\usepackage{mathrsfs} 

\usepackage{multicol}

\usepackage{mathrsfs}

\usepackage{cite}

\usepackage{cmap}

\usepackage{qrcode}

\allowdisplaybreaks[4]

\usepackage[nodayofweek]{datetime}

\usepackage[thinlines]{easybmat}



\renewcommand{\JournalYear}{2018}

\renewcommand{\JournalVolume}{22}

\renewcommand{\JournalNumber}{4}



\newdate{JournalDateSign}{29}{12}{2018}% Дата подписания Вестника в печать

\renewcommand{\JournalDateSign}{\displaydate{JournalDateSign}}

\newdate{JournalDatePrint}{16}{01}{2019}% Дата выхода Вестника из печати

\renewcommand{\JournalDatePrint}{\displaydate{JournalDatePrint}}

%\renewcommand{\JournalUPL}{16.56} % Усл. печ. л.

%\renewcommand{\JournalUIL}{16.53} % Уч.-изд. л.

\renewcommand{\JournalUPL}{15.24} % Усл. печ. л.

\renewcommand{\JournalUIL}{15.21} % Уч.-изд. л.

\renewcommand{\JournalRegNumber}{220/18} % Регистрационный номер

\renewcommand{\JournalOrderNumber}{2} % Номер заказа



%\renewcommand{\JournalRMonth}{Март} % Месяц выхода Вестника

%\renewcommand{\JournalEMonth}{March} % Месяц выхода Вестника

%\renewcommand{\JournalRMonth}{Июнь} % Месяц выхода Вестника

%\renewcommand{\JournalEMonth}{June} % Месяц выхода Вестника

%\renewcommand{\JournalRMonth}{Сентябрь} % Месяц выхода Вестника

%\renewcommand{\JournalEMonth}{September} % Месяц выхода Вестника

\renewcommand{\JournalRMonth}{Декабрь} % Месяц выхода Вестника

\renewcommand{\JournalEMonth}{December} % Месяц выхода Вестника



\newcommand{\totop}[1]{\raisebox{6pt}[1pt][2pt]{#1}}

\newcommand{\tottop}[1]{\raisebox{12pt}[1pt][2pt]{#1}} 

\newcommand{\totttop}[1]{\raisebox{18pt}[1pt][2pt]{#1}} 

\newcommand{\tottttop}[1]{\raisebox{24pt}[1pt][2pt]{#1}} 

\newcommand{\totttttop}[1]{\raisebox{30pt}[1pt][2pt]{#1}} 

\newcommand{\tobot}[1]{\raisebox{-6pt}[1pt][2pt]{#1}} 

\newcommand{\tobbot}[1]{\raisebox{-12pt}[1pt][2pt]{#1}} 

\newcommand{\tobbbot}[1]{\raisebox{-18pt}[1pt][2pt]{#1}}



\graphicspath{{./00PIC/}}

%

%

%%\samgtupubl % для генерации pdf, отдаваемого в типографию СамГТУ

%\pdfcrop % обрезка pdf для размещения на math-net.ru

%\draft % На корректуре

%%\filllastpage % Последняя страница не заполнена не полностью

%%\briefcomm % Статья является кратким сообщением

%

\vsgtu{1534}

%

\FirstPage{221}

\LastPage{235}

%%

%%

%\begin{document}









\renewcommand{\baselinestretch}{0.93}



\hfill \qrcode[hyperlink,height=.8in]{http://doi.org/10.14498/vsgtu1534}

\vspace{-.8in}







\udc{519.63}%Обработка текста. Подготовка текстов : Языки программирования. Компьютерные языки

\msc{65M06, 65M12}% Коды MSC см. http://www.ams.org/msc/



\rutitle{Об одном разностном методе решения задачи Трикоми для уравнения Лаврентьева--Бицадзе}



\rutitleshort{Об одном разностном методе решения задачи Трикоми\dots}



\rutitlehack{Об одном разностном методе решения \\ задачи~Трикоми~для~уравнения \\ Лаврентьева--Бицадзе}



\entitle{Finite-difference method for solving Tricomi problem for the Lavrent'ev--Bitsadze equation}

\entitleshort{Finite-difference method for solving Tricomi problem\dots}



\ruauthor{Ж.~А.~Балкизов, А.~А.~Сокуров}{Б\,а\,л\,к\,и\,з\,о\,в~Ж.~А., С\,о\,к\,у\,р\,о\,в~А.~А.}

\enauthor{G.~A.~Balkizov, A.~A.~Sokurov}{B\,a\,l\,k\,i\,z\,o\,v~G.~A., S\,o\,k\,u\,r\,o\,v~A.~A.}



\ruaffil{\ruaffid{4094}{Институт прикладной математики и автоматизации,\\

Россия, 360000, Нальчик, ул. Шортанова, 89~a}.}

\enaffil{\enaffid{4094}{Institute of Applied Mathematics and Automation, \\

89 a, Shortanov st., Nalchik, 360000,  Russian Federation}.}



 





\ruauthorsdetails{2}{

\fullauthor{\rupid{41451}{Жираслан Анатольевич Балкизов}} 

\CAuthor % Автор-корреспондент

\orcid{0000-0001-5329-7766} % ORCID ID автора 

\regalia{кандидат физико-математических наук} 

\position{заведующий отделом} 

\dept{отдел уравнений смешанного типа} 

\email{giraslan@yandex.ru}



\fullauthor{\rupid{124072}{Аслан Артурович Сокуров}} 

\orcid{0000-0001-9886-3602} % ORCID ID автора 

\position[n]{младший научный сотрудник} 

\dept{отдел теоретической и математической физики} 

\email[n]{asokuroff@gmail.com}

}







\enauthorsdetails{2}{

\fullauthor{\enpid{41451}{Giraslan A. Balkizov}} 

\CAuthor % Автор-корреспондент

\orcid{0000-0001-5329-7766} % ORCID ID автора 

\regalia{Cand. Phys. \& Math. Sci.} 

\position{Head of Department} 

\dept{Dept. of Mixed Type Equations} 

\email[n]{giraslan@yandex.ru}



\fullauthor{\enpid{124072}{Aslan A. Sokurov}} 

\orcid{0000-0001-9886-3602} % ORCID ID автора 

\position[n]{Junior Research Scientist} 

\dept{Dept. of Theoretical and Mathematical Physics} 

\email[n]{asokuroff@gmail.com}

}





\rukeywords{уравнение смешанного типа, задача Трикоми, априорная оценка, разностная схема, порядок аппроксимации, метод энергетических неравенств}



\enkeywords{equation of mixed type, Tricomi problem, a priori estimate, difference scheme, order of approximation, method of energy inequalities}



\ruabstract{Получена априорная оценка решения задачи Трикоми для уравнения

Лаврентьева--Бицадзе, из которой следует единственность регулярного решения. 

Предложен разностный метод нахождения приближенного решения исследуемой задачи. Получена априорная оценка решения разностной задачи, из которой следует сходимость построенной разностной схемы со вторым порядком точности.}



\enabstract{In this paper we obtain an a priori estimate for solution of Tricomi problem for 

the Lavrent'ev--Bitsadze equation, from which the uniqueness of regular solution follows. Presented a numerical finite-difference method for solving the investigated problem. We obtain an a priori estimate for solution of the difference scheme, from which follows the second-order convergence.}



\newdate{balkizova}{20}{03}{2017}

\date{\displaydate{balkizova}}

\newdate{balkizovb}{07}{05}{2017}

\revisiondate{\displaydate{balkizovb}}

\newdate{balkizovc}{12}{06}{2017}

\accepteddate{\displaydate{balkizovc}}



\newdate{balkizove}{10}{07}{2017}

\renewcommand{\JournalDataOnline}{\displaydate{balkizove}}





\makerutitle



%\def\sgn{\mathop{\rm sgn}\nolimits}





%Каждый пункт статьи должен начинаться с команды Название пункта

%после названия точку ставить не нужно. Пункты автоматически нумеруются.

%Если пункт нумеровать не нужно, то следует использовать в команде

% \Section необязательный параметр [n]:

%   \Section[n]{Имя пункта}

% Введение и заключение обычно не нумеруются.



\Section[n]{Введение}

Уравнение Лаврентьева"--~Бицадзе (ЛБ) в работах \cite{balk:c1, balk:c2}  было предложено как более простая модель уравнений смешанного типа в~смысле постановки и исследования смешанных задач. 

Методом сингулярных интегральных уравнений в~работах \cite{balk:c1, balk:c2, balk:c3, balk:c4, balk:c5, balk:c6} решена задача Трикоми для уравнения ЛБ для различных областей в~эллиптической части.

Исследованию однозначной разрешимости локальных и нелокальных краевых задач для уравнения ЛБ, а~также его обобщениям посвящено достаточно много работ. 

Отметим, например, работы \cite{balk:c7, balk:c8, balk:c9, balk:c10, balk:c11, balk:c12, balk:c13}.

Однако, несмотря на обширную библиографию работ по исследованию однозначной разрешимости различных форм дифференциальных задач, приближенным методам исследования краевых задач для уравнений смешанного типа уделено мало внимания. 

Среди них отметим работы \cite{balk:c14, balk:c15, balk:c16, balk:c17, balk:c18, balk:c19}, в~которых различными методами были получены приближенные решения задачи Трикоми для уравнения ЛБ с~первым порядком аппроксимации. 

В~данной работе получена априорная оценка решения задачи Трикоми для уравнения ЛБ в~случае, когда область эллиптичности ограничена прямоугольником, а~также предложен метод приближенного решения исследуемой задачи со вторым порядком точности. 

Отметим, что уравнение ЛБ  встречается при решении таких важных вопросов прикладного характера, как задачи теории бесконечно малых изгибаний поверхностей вращения, задачи безмоментной теории оболочек и~т.~д. 

Значительную роль такие уравнения играют и в задачах газовой динамики \cite{balk:c20}. 



\smallskip



\Section{Постановка задачи}

На евклидовой плоскости независимых переменных $x$ и $t$ рассмотрим уравнение Лаврентьева"--~Бицадзе

\begin{equation}

\label{lb1}

u_{xx}+\mathop{\rm sgn}(t)u_{tt}=-f(x, t), 

\end{equation}

где $\mathop{\rm sgn}(t)$ "--- знак числа $t$; $f(x, t)$ "--- заданная функция; $u=u(x,t)$ "--- искомая функция.



%\begin{figure}[h]

%\center{\includegraphics[width=6cm]{sss.eps}}

%\caption{Область $\Omega$.}

%\includegraphics[width=4.9cm]{sss.eps}

%\end{figure}







Через $\Omega$ обозначим область, ограниченную при $t<0$ характеристиками $AC: x+t=0$ и $CB: x-t=l$ уравнения \eqref{lb1}, выходящими из точек \linebreak $A\hm =(0, 0)$ и $B=(l, 0)$ 

и~пересекающимися в точке $C=(l/2, -l/2)$, а~при ${t>0}$ "--- прямоугольником с вершинами в точках $A$, $B$, $A_0=(0,\,T)$, ${B_0=(l,\,T)}$, $l>0$, $T>0$.

Верхнюю часть области $\Omega$ при $t>0$ обозначим через $\Omega^{+}$, а нижнюю часть при $t<0$ через $\Omega^{-}$.



Уравнение \eqref{lb1} является уравнением смешанного типа:  оно эллиптично в~области $\Omega^{+}$ и~гиперболично в $\Omega^{-}$. 



{\it Регулярным в области $\Omega$ решением} уравнения \eqref{lb1} назовем всякую функцию $u=u(x,t)$ из класса $C(\overline{\Omega})\cap C^1(\Omega)\cap C^2(\overline{\Omega}^{+})\cap C^2(\Omega^{-})$, при подстановке которой уравнение \eqref{lb1} обращается в тождество.



В работе рассматривается задача Трикоми для уравнения \eqref{lb1} в следующей постановке. 

%[умб 133].



\hypertarget{taskT}{}



{\small\sc Задача T.} {\it Найти регулярное в области $\Omega$ решение $u=u(x,t)$ уравнения \eqref{lb1}$,$ удовлетворяющее краевым условиям

\begin{gather}

\label{lb2}

u(l, t)=u_1 (t),\quad u(0, t)=u_3 (t),\quad 0 \leqslant t\leqslant T,

\\

\label{lb3}

u(x, T)=u_2 (x),\quad 0 \le x \le l,

\\

\label{lb4}

u(x,-x)=\psi (x),\quad 0 \le x\leqslant l/2,

\end{gather}

где $u_1(t),$ $u_2(x),$ $u_3(t),$ $\psi(x)$ "--- заданные достаточно гладкие функции$,$ для которых

 выполнены следующие условия  согласования}\/: 

$$u_3(0)=\psi(0),\quad u_3(T)=u_2(0),\quad  u_2(l)=u_1(T).$$





\smallskip



В работах \cite{balk:c14, balk:c15, balk:c16, balk:c17, balk:c18, balk:c19} были предложены различные методы нахождения приближенного решения задачи~\hyperlink{taskT}{T} для однородного уравнения Лаврентьева"--~Бицадзе с~первым порядком точности. 

В~настоящей работе предлагается конечно"=разностный метод, позволяющий найти приближенное решение задачи~\hyperlink{taskT}{T} для неоднородного уравнения Лаврентьева"--~Бицадзе со вторым порядком точности.



В дальнейшем будем предполагать, что $f(x,t) \in C(\overline{\Omega})$ и решение задачи \eqref{lb1}--\eqref{lb4} существует. 

Решение задачи Коши для уравнения \eqref{lb1} в области $\overline{\Omega}^{\,-}$ можно представить в~виде

\begin{equation}

\label{lb5}

u(x, t)=\frac{\tau(x+t)+\tau(x-t)}{2}+\frac{1}{2} \int _{x-t}^{x+t} \nu\left(\xi\right) d \xi+

\frac{1}{2} \int _{0}^{t}\int _{x-t+\eta}^{x+t-\eta} f (\xi,\eta) d\xi d\eta,

\end{equation}

где $\tau(x)=u(x, 0)$, $\nu(x)=u_t(x, 0).$

Учитывая условие \eqref{lb4}, из \eqref{lb5} получим

$$

u(x,-x)=\frac{\tau(0)+\tau\left(2x\right)}{2}+\frac{1}{2}\int _{2x}^{0}\nu\left(\xi\right) d \xi+

\frac{1}{2}\int _{0}^{-x}\int _{2x+\eta}^{-\eta}f (\xi,\eta) d\xi d\eta=\psi(x),

$$

или же

\begin{equation}

\label{lb6}

\frac{\tau(0)+\tau(x)}{2}+\frac{1}{2}\int _{x}^{0}\nu\left(\xi\right) d \xi+

\frac{1}{2}\int _{0}^{-x/2}\int _{x+\eta}^{-\eta}f (\xi,\eta) d\xi d\eta=\psi\Bigl(\frac{x}{2}\Bigr).

\end{equation}

Дифференцируя \eqref{lb6}, находим

\begin{equation}

\label{lb77}

\tau'(x)-\nu(x)=u_x (x, 0)-u_t (x, 0)=\psi'\Bigl(\frac{x}{2}\Bigr)+\int _{0}^{-x/2}f \left(x+\eta,\eta\right)d\eta,

\end{equation}

откуда

\begin{equation}

\label{lb7}

\nu(x)=\tau'(x)-\psi'\Bigl(\frac{x}{2}\Bigr)-\int _{0}^{-x/2}f (x+\eta,\eta)d\eta.

\end{equation}



Соотношение \eqref{lb7} "--- фундаментальное соотношение между функциями $\tau(x)$ и $\nu(x)$, принесенное из области $\Omega^{-}$ на линию $y=0$. 

Подставляя $\nu(x)$ из \eqref{lb7} в~\eqref{lb5}, находим решение первой задачи Дарбу для уравнения \eqref{lb1} в~области~$\Omega^{-}$:

\begin{multline}

u(x, t)=\tau(x+t)-\psi\Bigl(\frac{x+t}{2}\Bigr)+\psi\Bigl(\frac{x-t}{2}\Bigr)-

\\

\label{lb8}

-\frac{1}{2}\int _{x-t}^{x+t}\int _{0}^{-\xi/2}f \left(\xi+\eta,\eta\right)d\eta d\xi +\frac{1}{2} \int _{0}^{t}\int _{x-t+\eta}^{x+t-\eta} f (\xi,\eta) d\xi d\eta.

\end{multline}



Для определения значения искомой функции $u$ в области $\Omega^{+}$ с~учетом найденного  выше фундаментального соотношения приходим к~задаче \eqref{lb2}, \eqref{lb3} и \eqref{lb77} для уравнения \eqref{lb1}. 



%Если из них найдем $u$ в $\Omega^{+}$, то по значениям $\tau(x)=u(x, 0)$ найдем $u$ и в 

%$\Omega^{-}$ согласно \eqref{lb8}.



\smallskip



\Section{Априорная оценка} 

Справедлива следующая теорема.

\begin{theorem}[1]Для любого регулярного решения задачи \eqref{lb1}--\eqref{lb3} и~\eqref{lb77} с~однородными условиями \eqref{lb2}--\eqref{lb3} имеет место априорная оценка

\[

\left\Vert u \right\Vert_{0}^{2} \leqslant M_1 \left\Vert f \right\Vert_{0}^{2}+M_2 \left\Vert \chi \right\Vert_{0}^{2},

\]

где 

$$\chi(x)=\psi'\Bigl(\frac{x}{2}\Bigr)+\int _{0}^{-x/2}f \left(x+\eta,\eta\right)d\eta,$$

$M_1,$ $M_2$ "--- известные положительные числа$.$

\end{theorem}



\begin{proof}

Полагая условия \eqref{lb2},  \eqref{lb3} однородными, умножим уравнение \eqref{lb1} на $u(x, t)$ и~проинтегрируем по области $\Omega^{+}$:

\begin{multline}

\iint _{\Omega^{+}} u(x, t) \left[u_{xx}(x, t)+ u_{tt}(x, t)\right]dx dt=

\\

=\int _\Gamma -u(x, t) u_{t}(x, t)dx+u(x, t) u_{x}(x, t) dt-\iint _{\Omega^{+}} \left[u_{x}^{2}(x, t)+u_{t}^{2}(x, t)\right] dx dt=

\\

\label{lb9}

=-\iint _{\Omega^{+}}  u(x, t) f(x, t)dx dt,

\end{multline}

откуда

\begin{equation}

\label{lb10}

\int _{0}^{l} u(x, 0) u_{t}(x, 0)  dx+

\left\Vert u_{x} \right\Vert_{0}^{2}+\left\Vert u_{t} \right\Vert_{0}^{2}=\iint _{\Omega^{+}}  u(x, t) f(x, t)dx dt.

\end{equation}

Подставляя $u_{t}(x, 0)$ из \eqref{lb77} в \eqref{lb10}, приходим к равенству

\begin{equation}

\label{lb11}

-\int _{0}^{l} u(x, 0) \chi(x) dx+\left\Vert u_{x} \right\Vert_{0}^{2}+\left\Vert u_{t} \right\Vert_{0}^{2}=\iint _{\Omega^{+}}  u(x, t) f(x, t)dx dt.

\end{equation}





Пользуясь $\varepsilon$"=неравенством, убеждаемся в~справедливости неравенств

\begin{multline}

\iint _{\Omega^{+}} u(x, t) f(x, t)dx dt \leqslant 

\iint _{\Omega^{+}} \Bigl[\varepsilon_1 u^2 (x, t)+\frac{1}{4\varepsilon_1}f^{2}(x, t)\Bigr]dx dt=

\\

\label{lb12}

=\varepsilon_1 \left\Vert u \right\Vert_{0}^{2}+\frac{1}{4\varepsilon_1} \left\Vert f \right\Vert_{0}^{2}

\end{multline}

и

\begin{multline}

-\int _{0}^{l} u(x, 0) \chi(x)dx \geqslant 

\int _{0}^{l} \Bigl[-\varepsilon_2 u^2 (x, 0)-\frac{1}{4\varepsilon_2} \chi^2 (x)\Bigr] dx=

\\

\label{lb13}

=-\varepsilon_2 \int _{0}^{l} u^2 (x, 0)dx-\frac{1}{4\varepsilon_2} \left\Vert \chi \right\Vert_{0}^{2}.

\end{multline}

С учетом \eqref{lb12},  \eqref{lb13} неравенство \eqref{lb11} перепишется в следующем виде:

\begin{equation}

\label{lb14}

-\varepsilon_2 \int _{0}^{l} u^2 (x, 0)dx-\frac{1}{4\varepsilon_2} 

\left\Vert \chi \right\Vert_{0}^{2}+\left\Vert u_{x} \right\Vert_{0}^{2}+\left\Vert u_{t} \right\Vert_{0}^{2} \leqslant

\varepsilon_1 \left\Vert u \right\Vert_{0}^{2}+\frac{1}{4\varepsilon_1} \left\Vert f \right\Vert_{0}^{2}.

\end{equation}



Найдем теперь оценку для $\displaystyle \int _{0}^{l} u^2 (x, 0)dx$.

Воспользовавшись неравенством Коши"--~Буняковского с учетом  условия \eqref{lb3} получаем  

\begin{multline}

\int _{0}^{l} u^2 (x, 0)dx=\int _{0}^{l}\left(\int _{0}^{T} 

u_{t}(x, t)dt\right)^2 dx \leqslant \int _{0}^{l}\left(T \int _{0}^{T}u_{t}^{2}(x, t)dt\right)dx=

\\

\label{lb15}

=T \int _{0}^{l} \int _{0}^{T} u_{t}^{2}(x, t)dx dt=

T \iint _{\Omega^{+}}u_{t}^{2}(x, t)dx dt=T \left\Vert u_{t} \right\Vert_{0}^{2}.

\end{multline}



С учетом \eqref{lb15} неравенство \eqref{lb14} перепишется в следующем виде:

\begin{equation}

\label{lb16}

\left\Vert u_{x} \right\Vert_{0}^{2}+\left(1-\varepsilon_2 T\right) \left\Vert u_{t} \right\Vert_{0}^{2}\leqslant \varepsilon_1 \left\Vert u \right\Vert_{0}^{2}+\frac{1}{4\varepsilon_1} \left\Vert f \right\Vert_{0}^{2}+\frac{1}{4\varepsilon_2} \left\Vert \chi \right\Vert_{0}^{2}.

\end{equation}

Оценим далее $\left\Vert u_{x} \right\Vert_{0}^{2}$. Для этого заметим, что

\begin{equation}

\label{lb17}

u^2 (x, t)=\left(\int _{0}^{x} u_s \left(s,t\right)ds\right)^2 \leqslant 

x \int _{0}^{x} u_{s}^{2} \left(s,t\right)ds \leqslant x \int _{0}^{l} u_{x}^{2} (x, t)dx.

\end{equation}

Проинтегрируем неравенство \eqref{lb17} сначала по $x$ от 0 до $l$, а затем по $t$ от 0 до~$T$:

\begin{equation}

\label{lb18}

\int _{0}^{l} u^2 (x, t)dx \leqslant \frac{l^2}{2} \int _{0}^{l} u_{x}^{2} (x, t)dx;

\end{equation}

\begin{multline}

\int _{0}^{T}\left(\int _{0}^{l} u^2(x, t)dx\right) dt=\iint _{\Omega^{+}}u^{2}(x, t)dx dt \leqslant

\\

\label{lb19}

\leqslant \frac{l^2}{2} \int _{0}^{T}\left(\int _{0}^{l} u_{x}^{2}(x, t)dx\right) dt= \frac{l^2}{2} \iint _{\Omega^{+}} u_{x}^{2}(x, t)dx dt.

\end{multline}

Или окончательно

\begin{equation}

\label{lb20}

\frac{2}{l^2} \left\Vert u \right\Vert_{0}^{2} \leqslant \left\Vert u_x \right\Vert_{0}^{2}.

\end{equation}

Для оценки $\left\Vert u_{t} \right\Vert_{0}^{2}$ проинтегрируем неравенство 

$$

u^2 (x, t)=\left(-\int _{t}^{T} u_s (x, s)ds\right)^2 \le t \int _{0}^{T} u_{t}^{2} (x, t)dt

$$  

по области $\Omega^{+}$ и, рассуждая аналогично, найдем

\begin{equation}

\label{lb21}

\frac{2}{T^2} \left\Vert u \right\Vert_{0}^{2} \leqslant \left\Vert u_t \right\Vert_{0}^{2}.

\end{equation}

С учетом  \eqref{lb20} и \eqref{lb21} из \eqref{lb16} получаем

\begin{equation}

\label{lb22}

M \left\Vert u \right\Vert_{0}^{2} \leqslant \frac{1}{4\varepsilon_1} \left\Vert f \right\Vert_{0}^{2}+\frac{1}{4\varepsilon_2} \left\Vert \chi \right\Vert_{0}^{2},

\end{equation}  

где $$M=\frac{2}{l^2}+\frac{2}{T^2}-\frac{2 \varepsilon_2}{T}-\varepsilon_1.$$



В силу произвольности $\varepsilon_1$ и $\varepsilon_2$ число $M$ можно выбрать положительным. Тогда окончательно получим

\begin{equation}

\label{lb23}

\left\Vert u \right\Vert_{0}^{2} \leqslant M_1 \left\Vert f \right\Vert_{0}^{2}+M_2 \left\Vert \chi \right\Vert_{0}^{2},

\end{equation}

где $$M_1=\frac{1}{4\varepsilon_1 M},\quad M_2=\frac{1}{\varepsilon_2 M}.$$



Из априорной оценки \eqref{lb23} следует единственность регулярного решения исследуемой задачи.

\end{proof}







\smallskip



\Section{Разностная схема}

Для нахождения приближенного численного решения задачи \eqref{lb1}--\eqref{lb3}, \eqref{lb77} для уравнения \eqref{lb1} в~области $\overline{\Omega}^{+}$ воспользуемся методом конечных разностей.

В~$\overline{\Omega}^{+}$ введем сетку 

$$\overline{\omega}_{h_x h_t}=\overline{\omega}_{h_x}  \times \overline{\omega}_{h_t},$$ 

где 

$\overline{\omega}_{h_x}=\{x_i=i h_x,\,i=\overline{0,\,N_x}\}$, 

$\overline{\omega}_{h_t}=\{t_j=j h_t,\,j=\overline{0,\,N_t}\}$, 

$h_x=l/N_x$, $h_t=T/N_t$, $N_x$ и $N_t$ "--- натуральные числа. 



%\begin{figure}[h]

%\center{\includegraphics[width=6cm]{ssss.eps}}

%\caption{Область $\Omega$.}

%\includegraphics[width=4.9cm]{sss.eps}

%\end{figure}

%Найдем сначала разностную аппроксимация для условия $\ref{lb7}$. Справедлива следующая 

%{\bf Лемма}

%Условию $\ref{lb77}$ соответствует разностная аппроксимация вида 

%\begin{equation}

%\label{lb277}

%\left(y_{\mathring{x}}\right)_{i}^{0}-\left(y_{t}\right)_{i}^{0}-

%\frac{h_t}{2} \left(y_{\overline{x}x}\right)_{i}^{0}=w_i,\qquad i=\overline{1,N_x-1},

%\end{equation}

%где $\varphi_{i}^{j}=f(x_i,t_j),\,w_i=\frac{h_t}{2} f^{+}(x_i,0)+\psi'\left(\frac{x_i}{2}\right)

%+\int _{0}^{\frac{-x_i}{2}}f^{-}(x_i+\eta,\eta)d\eta$, порядок аппроксимации которого равен двум.

%Действительно,  



Исходной дифференциальной задаче поставим в соответствие следующую разностную схему:

\begin{gather}

\label{lb24}

\left(y_{\overline{x}x}\right)_{i}^{j}+\left(y_{\overline{t}t}\right)_{i}^{j}=-\varphi_{i}^{j},\quad

i=\overline{1,\,N_x-1},\,j=\overline{1,\,N_t-1},

\\

\label{lb25}

y_{N_x}^{j}=u_1(t_j),\quad y_{0}^{j}=u_3(t_j),\quad j=\overline{0,\,N_t},

\\

\label{lb26}

y_{i}^{N_t}=u_2(x_i),\quad i=\overline{1,\,N_x-1},

\\

\label{lb27}

\left(y_{\mathring{x}}\right)_{i}^{0}-\left(y_{t}\right)_{i}^{0}-

\frac{h_t}{2} \left(y_{\overline{x}x}\right)_{i}^{0}=w_i,\quad i=\overline{1,\,N_x-1}.

\end{gather}

где 

$$\varphi_{i}^{j}=f(x_i,t_j),\quad  w_i=\frac{h_t}{2} f(x_i,0)+\psi'\Bigl(\frac{x_i}{2}\Bigr)

+\int _{0}^{{-x_i}/{2}}f(x_i+\eta,\eta)d\eta.$$



Введем следующие  обозначения:

\[

\left\Vert y\right.]|_{0}^{2}=\sum\limits_{j=1}^{N_t}\sum\limits_{i=1}^{N_x}\left(y_{i}^{j}\right)^2 h_t h_x,\quad 

\left\Vert y \right\Vert_{C(\overline{\omega})}=\max\limits_{x\in\,\overline{\omega}}|y(x)|,

\quad 

(y,v)=\sum\limits_{j=1}^{N_t-1}\sum\limits_{i=1}^{N_x-1}y_{i}^{j} v_{i}^{j} h_t h_x.

\]





\begin{theorem}[2]Для любого решения $y(x),$ $x \in \overline{\omega}_{h_x h_t},$ задачи  

\eqref{lb24}--\eqref{lb27} с~однородными условиями \eqref{lb25}--\eqref{lb26} справедлива априорная оценка

\[

\left\Vert y\right.]|_{0}^{2}\leqslant M_3\left\Vert w \right.]|_{0}^{2}+M_4 \left\Vert \varphi \right.]|_{0}^{2},

\]

где $M_3,$ $M_4$ "--- известные положительные числа$.$

\end{theorem}



\begin{proof}

Пользуясь методом энергетических неравенств \cite[c.~110]{balk:c21}, умножим \eqref{lb24} скалярно на $y$:

\begin{equation}

\label{lb28}

\left(y,y_{\overline{x}x}\right)+\left(y,y_{\overline{t}t}\right)=-(y,\varphi).

\end{equation}

Проведем следующие вычисления:

\begin{multline}

\left(y,y_{\overline{x}x}\right)=\sum\limits_{j=1}^{N_t-1}\sum\limits_{i=1}^{N_x-1}

\left(y_{\overline{x}x}\right)_{i}^{j} y_{i}^{j} h_t h_x=

\sum\limits_{j=1}^{N_t-1}\sum\limits_{i=1}^{N_x-1} \frac{(y_{\overline{x}})_{i+1}^{j}-(y_{\overline{x}})_{i}^{j}}{h_x} y_{i}^{j} h_x h_t=

\\

=\sum\limits_{j=1}^{N_t-1}\sum\limits_{i=1}^{N_x-1}\left[(y_{\overline{x}})_{i+1}^{j} y_{i}^{j}

-(y_{\overline{x}})_{i}^{j} y_{i-1}^{j}-(y_{\overline{x}})_{i}^{j}  y_{i}^{j}

+(y_{\overline{x}})_{i}^{j}y_{i-1}^{j}\right] h_t=

\\

=\sum\limits_{j=1}^{N_t-1}\sum\limits_{i=1}^{N_x-1}\left[(y_{\overline{x}})_{i+1}^{j} y_{i}^{j}

-(y_{\overline{x}})_{i}^{j} y_{i-1}^{j}\right]h_t-\sum\limits_{j=1}^{N_t-1}\sum\limits_{i=1}^{N_x-1}

\frac{y_{i}^{j}-y_{i-1}^{j}}{h_x}(y_{\overline{x}})_{i}^{j} h_t h_x=

\\

=\sum\limits_{j=1}^{N_t-1}\left[(y_{\overline{x}})_{N_x}^{j} y_{N_x-1}^{j}

-(y_{\overline{x}})_{1}^{j} y_{0}^{j}\right]h_t-\sum\limits_{j=1}^{N_t-1}\sum\limits_{i=1}^{N_x-1}

\left[(y_{\overline{x}})_{i}^{j} \right]^2 h_t h_x=

\\

=-\sum\limits_{j=1}^{N_t-1}\frac{y_{N_x}^{j}- y_{N_x-1}^{j}}{h_x}(y_{\overline{x}})_{N_x}^{j} h_x h_t-\sum\limits_{j=1}^{N_t-1}\sum\limits_{i=1}^{N_x-1}\left[(y_{\overline{x}})_{i}^{j} \right]^2 h_t h_x=

\\

=-\sum\limits_{j=1}^{N_t-1}\left[\left(y_{\overline{x}}\right)_{N_x}^{j}\right]^2 h_t h_x-\sum\limits_{j=1}^{N_t-1}\sum\limits_{i=1}^{N_x-1}\left[(y_{\overline{x}})_{i}^{j} \right]^2 h_t h_x=-\sum\limits_{j=1}^{N_t-1}\sum\limits_{i=1}^{N_x}\left[(y_{\overline{x}})_{i}^{j} \right]^2 h_t h_x=

\\

=-\sum\limits_{j=1}^{N_t-1}\sum\limits_{i=1}^{N_x}\left[(y_{\overline{x}})_{i}^{j} \right]^2 h_t h_x-\sum\limits_{i=1}^{N_x}\left[(y_{\overline{x}})_{i}^{N_t} \right]^2 h_t h_x 

=

\\

=-\sum\limits_{j=1}^{N_t}\sum\limits_{i=1}^{N_x}\left[(y_{\overline{x}})_{i}^{j} \right]^2 h_t h_x=-\left\Vert y_{\overline{x}}\right.]|_{0}^{2}

\end{multline}

и

\begin{multline}

\left(y,y_{\overline{t}t}\right)=\sum\limits_{i=1}^{N_x-1}\sum\limits_{j=1}^{N_t-1}

\left(y_{\overline{t}t}\right)_{i}^{j} y_{i}^{j} h_t h_x=

\sum\limits_{i=1}^{N_x-1}\sum\limits_{j=1}^{N_t-1} \frac{(y_{\overline{t}})_{i}^{j+1}-(y_{\overline{t}})_{i}^{j}}{h_t} y_{i}^{j} h_x h_t=

\\

=\sum\limits_{i=1}^{N_x-1}\sum\limits_{j=1}^{N_t-1}\left[(y_{\overline{t}})_{i}^{j+1} y_{i}^{j}

-(y_{\overline{t}})_{i}^{j} y_{i}^{j-1}-(y_{\overline{t}})_{i}^{j}  y_{i}^{j}

+(y_{\overline{t}})_{i}^{j}y_{i}^{j-1}\right] h_x=

\\

=\sum\limits_{i=1}^{N_x-1}\sum\limits_{j=1}^{N_t-1}\left[(y_{\overline{t}})_{i}^{j+1} y_{i}^{j}

-(y_{\overline{t}})_{i}^{j} y_{i}^{j-1}\right]h_x-\sum\limits_{i=1}^{N_x-1}\sum\limits_{j=1}^{N_t-1}

\frac{y_{i}^{j}-y_{i}^{j-1}}{h_t}(y_{\overline{t}})_{i}^{j} h_t h_x=

\\

=\sum\limits_{i=1}^{N_x-1}\left[(y_{\overline{t}})_{i}^{N_t} y_{i}^{N_t-1}

-(y_{\overline{t}})_{i}^{1} y_{i}^{0}\right]h_x-\sum\limits_{i=1}^{N_x-1}\sum\limits_{j=1}^{N_t-1}

\left[(y_{\overline{t}})_{i}^{j} \right]^2 h_t h_x=

\\

=-\sum\limits_{i=1}^{N_x-1}\frac{y_{i}^{N_t}

- y_{i}^{N_t-1}}{h_t}(y_{\overline{t}})_{i}^{N_t} h_t h_x-\sum\limits_{i=1}^{N_x-1}(y_{\overline{t}})_{i}^{1} y_{i}^{0} h_x-\sum\limits_{i=1}^{N_x-1}\sum\limits_{j=1}^{N_t-1}

\left[(y_{\overline{t}})_{i}^{j} \right]^2 h_t h_x=

\\

=

-\sum\limits_{i=1}^{N_x-1}\left[\left(y_{\overline{t}}\right)_{i}^{N_t}\right]^2 h_t h_x-\sum\limits_{i=1}^{N_x-1}(y_t)_{i}^{0} y_{i}^{0} h_x-\sum\limits_{i=1}^{N_x-1}\sum\limits_{j=1}^{N_t-1}

\left[(y_{\overline{t}})_{i}^{j} \right]^2 h_t h_x=

\\

=-\sum\limits_{i=1}^{N_x-1}(y_t)_{i}^{0} y_{i}^{0} h_x-\sum\limits_{i=1}^{N_x-1}\sum\limits_{j=1}^{N_t}\left[(y_{\overline{t}})_{i}^{j} \right]^2 h_t h_x=

\\

=-\sum\limits_{i=1}^{N_x-1}(y_t)_{i}^{0} y_{i}^{0} h_x-\sum\limits_{i=1}^{N_x-1}\sum\limits_{j=1}^{N_t}\left[(y_{\overline{t}})_{i}^{j} \right]^2 h_t h_x-\sum\limits_{j=1}^{N_t}\left[(y_{\overline{t}})_{N_x}^{j} \right]^2 h_t h_x=

\\

=-\sum\limits_{i=1}^{N_x-1}(y_t)_{i}^{0} y_{i}^{0} h_x-\sum\limits_{i=1}^{N_x}\sum\limits_{j=1}^{N_t}\left[(y_{\overline{t}})_{i}^{j} \right]^2 h_t h_x=-\sum\limits_{i=1}^{N_x-1}(y_t)_{i}^{0} y_{i}^{0} h_x-\left\Vert y_{\overline{t}}\right.]|_{0}^{2}.

\end{multline}

Таким образом, 

\begin{equation}

\label{lb29}

\left\Vert y_{\overline{x}}\right.]|_{0}^{2}+\left\Vert y_{\overline{t}}\right.]|_{0}^{2}+\sum\limits_{i=1}^{N_x-1}(y_t)_{i}^{0} y_{i}^{0} h_x=(y,\varphi).

\end{equation}

Выразим $(y_t)_{i}^{0}$ из \eqref{lb27} и подставим в \eqref{lb29}. Тогда

\begin{equation}

\label{lb30}

\left\Vert y_{\overline{x}}\right.]|_{0}^{2}+\left\Vert y_{\overline{t}}\right.]|_{0}^{2}+\sum\limits_{i=1}^{N_x-1}\left[\left(y_{\mathring{x}}\right)_{i}^{0}-w_i-\frac{h_t}{2}

(y_{\overline{x}x})_{i}^{0}\right]y_{i}^{0} h_x=(y,\varphi).

\end{equation}

Легко заметить, что 

\begin{equation}

\label{lb31}

\sum\limits_{i=1}^{N_x-1}(y_{\mathring{x}})_{i}^{0} y_{i}^{0} h_x=y_{N_x}^{0}y_{N_x-1}^{0}-y_{0}^{0}y_{1}^{0}=0,

\end{equation}

\begin{multline}

\frac{h_t h_x}{2} \sum\limits_{i=1}^{N_x-1}(y_{\overline{x}x})_{i}^{0}y_{i}^{0}=

\frac{h_t}{2}(y_{N_x}^{0} \left.y_{\overline{x}}\right._{N_x}^{0}-y_{0}^{0}\left.y_{\overline{x}}\right._{1}^{0})-\frac{h_t h_x}{2} \sum\limits_{i=1}^{N_x}\left.y_{\overline{x}}^2\right._{i}^{0}=

\\

\label{lb32}

=- \frac{h_t h_x}{2} \sum\limits_{i=1}^{N_x}\left.y_{\overline{x}}^2\right._{i}^{0}.

\end{multline}

%$$

%=\left[(y y_{\overline{x}})_{N_x}^{0}-(y y_{\overline{x}})_{1}^{0}\right]\frac{h_t}{2}-

%\sum\limits_{i=1}^{N_x-1}\left.y_{\overline{x}}^2\right._{i}^{0} \frac{h_t h_x}{2}=

%$$

%\begin{equation}

%\label{lb32}

%=-y_{1}^{0}(y_{1}^{0}-y_{0}^{0})\frac{h_t}{2}-\sum\limits_{i=1}^{N_x-1}\left.y_{\overline{x}}^2\right._{i}^{0} %\frac{h_t h_x}{2}=-\left(\frac{\left.y_{1}^{0}\right.^2}{h_x}+\sum\limits_{i=1}^{N_x-1}\left.y_{\overline{x}}%^2\right._{i}^{0} \right)\frac{h_t h_x}{2}.

%\end{equation}

С учетом \eqref{lb31}, \eqref{lb32} из \eqref{lb30} получаем

\begin{equation}

\label{lb33}

\left\Vert y_{\overline{x}}\right.]|_{0}^{2}+\left\Vert y_{\overline{t}}\right.]|_{0}^{2}+\frac{h_t h_x}{2} \sum\limits_{i=1}^{N_x}\left.y_{\overline{x}}^2\right._{i}^{0}-\sum\limits_{i=1}^{N_x-1} w_i y_{i}^{0} h_x=(y,\varphi),

\end{equation}

откуда 

\begin{equation}

\label{lb34}

\left\Vert y_{\overline{x}}\right.]|_{0}^{2}+\left\Vert y_{\overline{t}}\right.]|_{0}^{2}-h_x\sum\limits_{i=1}^{N_x-1} w_i y_{i}^{0}\leqslant (y,\varphi).

\end{equation}

Применим к третьему слагаемому слева в \eqref{lb34} $\varepsilon$"=неравенство: 

\begin{equation}

\label{lb35}

h_x\sum\limits_{i=1}^{N_x-1} w_i y_{i}^{0} \leqslant

h_x \sum\limits_{i=1}^{N_x} \Bigl(\frac{1}{4\varepsilon_1} w_{i}^{2}+\varepsilon_1 \left.y_{i}^{0}\right.^2\Bigr).

\end{equation}

%$$

%\left\Vert y_{\overline{x}} \right\Vert^{2}+\left\Vert y_{\overline{t}} \right\Vert^{2}-

%\varepsilon_1 \sum\limits_{i=1}^{N_x-1}  w_{i}^{2} h_x-

%\frac{1}{4\varepsilon_1}\sum\limits_{i=1}^{N_x-1} \left.y_{i}^{0}\right.^2 h_x\leqslant

%$$

Тогда из \eqref{lb34} получим

\begin{gather}

\left\Vert y_{\overline{x}}\right.]|_{0}^{2}+\left\Vert y_{\overline{t}}\right.]|_{0}^{2}-

\frac{1}{4\varepsilon_1}\sum\limits_{i=1}^{N_x}  w_{i}^{2} h_x-

\varepsilon_1 \sum\limits_{i=1}^{N_x} \left.y_{i}^{0}\right.^2 h_x \leqslant\varepsilon_2 \left\Vert y \right.]|_{0}^{2}+\frac{1}{4\varepsilon_2} \left\Vert \varphi \right.]|_{0}^{2},

\\

\label{lb36}

\left\Vert y_{\overline{x}}\right.]|_{0}^{2}+\left\Vert y_{\overline{t}}\right.]|_{0}^{2}-\varepsilon_2 \left\Vert y \right.]|_{0}^{2}-\varepsilon_1\sum\limits_{i=1}^{N_x} \left.y_{i}^{0}\right.^2 h_x \leqslant \frac{1}{4\varepsilon_1}\left\Vert w \right.]|_{0}^{2}+\frac{1}{4\varepsilon_2} \left\Vert \varphi \right.]|_{0}^{2},

\end{gather}

где $\left\Vert w \right.]|_{0}^{2}=\sum _{i=1}^{N_x}  w_{i}^{2} h_x$.\\

Оценим $\sum _{i=1}^{N_x} \left.y_{i}^{0}\right.^2 h_x$. Для этого заметим, что 

$$

y_{i}^{0}=-\sum\limits_{j=1}^{N_t}(y_{\overline{t}})_{i}^{j} h_t,

$$ 

откуда

\begin{equation}

\label{lb37}

(y_{i}^{0})^2=\biggl(\sum\limits_{j=1}^{N_t}(y_{\overline{t}})_{i}^{j} h_t\biggr)^2\leqslant

T \sum\limits_{j=1}^{N_t}\left[(y_{\overline{t}})_{i}^{j}\right]^2 h_t.

\end{equation} 

Суммируя обе части неравенства \eqref{lb37} по $i$ от $i=1$ до $i=N_x$, приходим к~оценке 

\begin{equation}

\label{lb38}

\sum\limits_{i=1}^{N_x}(y_{i}^{0})^2 h_x \leqslant T\sum\limits_{i=1}^{N_x}\sum\limits_{j=1}^{N_t}\left[(y_{\overline{t}})_{i}^{j}\right]^2 h_t h_x=T\left\Vert y_{\overline{t}}\right.]|_{0}^{2}.

\end{equation}

С учетом \eqref{lb38} неравенство \eqref{lb36} примет вид

\begin{equation}

\label{lb39}

\left\Vert y_{\overline{x}}\right.]|_{0}^{2}+\left(1-\varepsilon_1 T \right)\left\Vert y_{\overline{t}}\right.]|_{0}^{2}-\varepsilon_2 \left\Vert y \right.]|_{0}^{2}\leqslant \frac{1}{4\varepsilon_1}\left\Vert w \right.]|_{0}^{2}+\frac{1}{4\varepsilon_2} \left\Vert \varphi \right.]|_{0}^{2}.

\end{equation}

Найдем теперь разностные аналоги неравенств \eqref{lb20} и \eqref{lb21}. 

Имеем  

$$ 

(y_{i}^{j})^2=\biggl(\sum\limits_{k=1}^{i}\left(y_{\overline{x}}\right)_{k}^{j} h_x\biggr)^2\leqslant

ih_x \sum\limits_{k=1}^{i}\left[\left(y_{\overline{x}}\right)_{k}^{j}\right]^2 h_x \leqslant

ih_x \sum\limits_{k=1}^{N_x}\left[\left(y_{\overline{x}}\right)_{k}^{j}\right]^2 h_x,

$$



$$

\sum\limits_{i=1}^{N_x} (y_{i}^{j})^2 h_x 

\leqslant \biggl(\sum\limits_{k=1}^{N_x}\left[\left(y_{\overline{x}}\right)_{k}^{j}\right]^2 h_x\biggr)\biggl(\sum\limits_{i=1}^{N_x}ih_{x}^{2}\biggr)=

\frac{l(l+h_x)}{2}\sum\limits_{i=1}^{N_x}\left[\left(y_{\overline{x}}\right)_{i}^{j}\right]^2 h_x,

$$



$$

\sum\limits_{j=1}^{N_t}\sum\limits_{i=1}^{N_x} (y_{i}^{j})^2 h_t h_x \leqslant

\frac{l(l+h_x)}{2}\sum\limits_{j=1}^{N_t}\sum\limits_{i=1}^{N_x}\left[\left(y_{\overline{x}}\right)_{i}^{j}\right]^2 h_t h_x,

$$

откуда

\begin{equation}

\label{lb40}

\frac{2}{l(l+h_x)}\left\Vert y\right.]|_{0}^{2}\leqslant \left\Vert y_{\overline{x}}\right.]|_{0}^{2}.

\end{equation}

Аналогично получаем оценку

\begin{equation}

\label{lb41}

\frac{2}{T^2}\left\Vert y\right.]|_{0}^{2}\leqslant \left\Vert y_{\overline{t}}\right.]|_{0}^{2}.

\end{equation}

С учетом \eqref{lb40},  \eqref{lb41} из \eqref{lb39} приходим к оценке вида

\begin{equation}

\label{lb42}

M\left\Vert y\right.]|_{0}^{2}\leqslant \frac{1}{4\varepsilon_1}\left\Vert w \right.]|_{0}^{2}+\frac{1}{4\varepsilon_2} \left\Vert \varphi \right.]|_{0}^{2},

\end{equation}

где $$M= \frac{2}{l(l+h_x)}+\frac{2}{T^2}-\frac{2\varepsilon_1}{T}-\varepsilon_2 .$$

Выбирая значения постоянных $\varepsilon_1$ и $\varepsilon_2$ так, чтобы $M>0$, из \eqref{lb42} окончательно приходим к априорной оценке

\begin{equation}

\label{lb43}

\left\Vert y\right.]|_{0}^{2}\leqslant M_3\left\Vert w \right.]|_{0}^{2}+M_4 \left\Vert \varphi \right.]|_{0}^{2},

\end{equation}

где $$M_3=\frac{1}{4\varepsilon_1 M},\quad M_4=\frac{1}{4\varepsilon_2 M}.$$



Из априорной оценки \eqref{lb43} следует существование и единственность решения разностной задачи 

\eqref{lb24}--\eqref{lb27}.

\end{proof}





\smallskip



\Section{Порядок аппроксимации и сходимость разностной схемы}

Пусть $u$ "--- точное решение задачи \eqref{lb1}--\eqref{lb3}, \eqref{lb77} из класса $C^4(\overline{\Omega}^{+})$ (завышенного условия гладкости на точное решение $u$ будем требовать для получения 2-го порядка аппроксимации), a $y$ "--- решение соответствующей разностной задачи \eqref{lb24}--\eqref{lb27}. 

Тогда для погрешности $z=y-u$ получаем задачу

\begin{gather}

\label{lb44}

\left(z_{\overline{x}x}\right)_{i}^{j}+\left(z_{\overline{t}t}\right)_{i}^{j}=-\overline{\varphi}_{i}^{j},\quad

i=\overline{1,\,N_x-1},\,j=\overline{1,\,N_t-1},

\\

\label{lb45}

z_{N_x}^{j}=0, \quad z_{0}^{j}=0,\quad j=\overline{0,\,N_t},

\\

\label{lb46}

z_{i}^{N_t}=0,\quad i=\overline{1,\,N_x-1},

\\

\label{lb47}

\left(z_{\mathring{x}}\right)_{i}^{0}-\left(z_{t}\right)_{i}^{0}-

\frac{h_t}{2} \left(z_{\overline{x}x}\right)_{i}^{0}=\overline{w}_i,\qquad i=\overline{1,\,N_x-1},

\end{gather}

где 

$$\overline{\varphi}_{i}^{j}=\left(u_{\overline{x}x}\right)_{i}^{j}+

\left(u_{\overline{t}t}\right)_{i}^{j}+\varphi_{i}^{j},

\quad 

\overline{w}_i=-\bigl(\left(u_{\mathring{x}}\right)_{i}^{0}-\left(u_{t}\right)_{i}^{0}-\frac{h_t}{2} \left(u_{\overline{x}x}\right)_{i}^{0}-w_i\bigr).

$$





Пользуясь формулой Тейлора, оценим порядок малости сеточных функций $\overline{\varphi}$ и $\overline{w}$. 

Для $\overline{\varphi}$ получаем

\begin{multline}

\overline{\varphi}_{i}^{j}=\left(u_{\overline{x}x}\right)_{i}^{j}+

\left (u_{\overline{t}t}\right)_{i}^{j}+\varphi_{i}^{j}=

\Bigl(\frac{\partial^2 u}{\partial x^2}\left(x_i,t_j\right)+\frac{\partial^2 u}{\partial t^2}\left(x_i,t_j\right)+f\left(x_i,t_j\right)\Bigr)+

\\

+\frac{h_x^2}{12}\frac{\partial^4 u}{\partial x^4}\left(\xi_i,t_j\right)+\frac{h_t^2}{12}\frac{\partial^4 u}{\partial t^4}\left(x_i,\tau_j\right)=O\left(h_x^2+h_t^2\right),

\end{multline}

где 

$$\xi_i=x_i+\theta_1 h_x,

\quad  

\left|\theta_1\right|\le 1,

\quad 

\tau_j=t_j+\theta_2 h_t,\quad  \left|\theta_2\right|\le 1.

$$



Далее найдем порядок аппроксимации граничного условия \eqref{lb77}.

Для этого воспользуемся следующими разложениями: 

\begin{gather*}

\left(u_{\mathring{x}}\right)_{i}^{0}=\frac{\partial u}{\partial x}\left(x_i,0\right)+\frac{h_x^2}{6}\frac{\partial^3 u}{\partial x^3} (\overline{\xi}_i, 0 ),

\\

\left(u_{t}\right)_{i}^{0}=\frac{\partial u}{\partial t}\left(x_i,0\right)+\frac{h_t}{2}\frac{\partial^2 u}{\partial t^2}\left(x_i,0\right)+\frac{h_t^2}{6}\frac{\partial^3 u}{\partial t^3} (x_i,\overline{\tau_i} ),

\\

\left(u_{\overline{x}x}\right)_{i}^{0}=\frac{\partial^2 u}{\partial x^2}\left(x_i,0\right)+

\frac{h_x^2}{12}\frac{\partial^4 u}{\partial x^4} (\overline{\overline{\xi}}_i,0 ),

\end{gather*}

где 

$\overline{\xi_i}$,

$\overline{\overline{\xi}}_i$ "--- некоторые промежуточные точки из интервала $(x_i-h_x,\,x_i+h_x)$, а~$\overline{\tau}_i \in (0,\,h_t)$. 

С~учетом приведенных выше разложений из \eqref{lb27} находим

\begin{multline}

\left(u_{\mathring{x}}\right)_{i}^{0}-\left(u_{t}\right)_{i}^{0}-

\frac{h_t}{2} \left(u_{\overline{x}x}\right)_{i}^{0}-w_i=

\\

=\biggl(\frac{\partial u}{\partial x}\left(x_i,0\right)-\frac{\partial u}{\partial t}\left(x_i,0\right)-\psi'\Bigl(\frac{x_i}{2}\Bigr)-

\int _{0}^{{-x_i}/{2}}f^{-}(x_i+\eta,\eta)d\eta \biggr)-

\\

-\frac{h_t}{2}

\Bigl(\frac{\partial^2 u}{\partial x^2}\left(x_i,0\right)+

\frac{\partial^2 u}{\partial t^2}\left(x_i,0\right)+f\left(x_i,0\right) \Bigr)+

\frac{h_x^2}{6}\frac{\partial^3 u}{\partial x^3}

(\overline{\xi}_i, 0 )-

\\

-\frac{h_t^2}{6}\frac{\partial^3 u}{\partial t^3}\left(x_i,\overline{\tau_i}\right)-\frac{h_t h_x^2}{12}\frac{\partial^4 u}{\partial x^4}

(\overline{\overline{\xi}}_i,0 )=O(h_x^2+h_t^2 ).

\end{multline} 

Таким образом,

$$

\overline{\varphi}=O\left(h_x^2+h_t^2\right),\qquad \overline{w}=O\left(h_x^2+h_t^2\right).

$$



Для $z$ справедлива априорная оценка \eqref{lb43}, откуда и будет следовать сходимость разностной схемы со скоростью $O\left(h_x^2+h_t^2\right)$ в норме $\left\Vert \cdot \right.]|_{0}$. 





\smallskip



\Section{Численные расчеты}

Численные расчеты с использованием разностной схемы \eqref{lb24}--\eqref{lb27} были проведены для тестовой задачи в~случае, когда входные данные заданы следующим образом:

\begin{gather}

f(x, t)=

\begin{cases}

(-2+36 t^2)\cos(6x)+(2-25x^2)\sin(5t), & t>0,\\

10t-2\cos(t)\cos(x)\cosh(t)-\cos(x)\sin(t)\sh(t),& t<0,

\end{cases}

\\

u_1(t)=-t^2 \cos(6)+\sin(5t),\quad u_2(x)=-\cos(6x)+x^2 \sin(5),

\\

u_3(t)=-t^2,\quad \psi(x)=-5x^3+\cos(x)\sin(x)\sh(x),

\\

l= 1,\quad T=1.

\end{gather}

Точное решение задачи

$$

u(x, t)=

\begin{cases}

x^2\sin(5t)-t^2\cos(6x), & t>0,\\

5x^2 t+\sin(t)\sh(t)\cos(x), & t<0.

\end{cases}$$





\begin{table}[h!]

\renewcommand{\tabcolsep}{.75mm}

\centering 

\begin{minipage}[t]{128mm} \small \centering { \footnotesize

Погрешность и порядок сходимости в нормах $\left\Vert \cdot \right.]|_0$ и $\left\Vert \cdot \right\Vert_{C(\overline{\omega})}$



[Error estimation and rate of convergence of the $\left\Vert \cdot \right.]|_0$ and $\left\Vert \cdot \right\Vert_{C(\overline{\omega})}$ norms] } \vspace{-3mm} 

\begin{tabular}{c||c|c|c|c}

\hline 

& \footnotesize Error estimation   

& \footnotesize Rate of convergence 

& \footnotesize Error estimation   

& \footnotesize Rate of convergence \\

$h_x=h_t$ 

& \footnotesize of the 

$\left\Vert z \right.]|_0$ norm

& \footnotesize of the $\left\Vert z \right.]|_0$ norm 

& \footnotesize of the $\left\Vert z \right\Vert_{C(\overline{\omega})}$ norm

& \footnotesize of the $\left\Vert z \right\Vert_{C(\overline{\omega})}$ norm\\	

\hline\hline

1/20 &  0.00172 &  &  0.00373 & \\

%\hline

1/40 &  $4.26217\cdot 10^{-4}$ &  2.01029 &   $9.22934\cdot 10^{-4}$ &  2.01678\\

%\hline

1/80 & $1.06206\cdot 10^{-4}$ & 2.00472 &  $2.30968\cdot 10^{-4}$ &  1.99853\\

%\hline

1/160 &  $2.65094\cdot 10^{-5}$ & 2.00228 & $5.77445\cdot 10^{-5}$ &  1.99993\\

\hline

\end{tabular}

\end{minipage}

\end{table}





Порядок сходимости вычислялся согласно формуле 

$$\log_{h_1/h_2}(\Delta_1/\Delta_2),$$

где $\Delta_i$ "--- норма погрешности, соответствующая шагу $h_i=h_x=h_t$.



В таблице приведены значения погрешности и вычисленный на их основе порядок сходимости в нормах $\left\Vert \cdot \right.]|_0$ и $\left\Vert \cdot \right\Vert_{C(\overline{\omega})}$. 

Как видно, сгущение сетки приводит к уменьшению погрешности и порядок сходимости стремится к~двум. 

Эти выводы согласуются с результатами, полученными при исследовании свойств разностной схемы (\ref{lb24})--(\ref{lb27}).





\AdditionalContent{Конкурирующие интересы}{Мы не имеем конкурирующих интересов.} 

\AdditionalContent{Авторский вклад и ответственность}{Все  авторы  принимали  участие  в  разработке  концепции  \mbox{статьи}  и  в  написании  рукописи. Авторы несут  полную  ответственность  за  предоставление  окончательной рукописи в печать. Окончательная  версия рукописи была одобрена всеми авторами.} 

\AdditionalContent{Финансирование}{Исследование выполнялось без финансирования.}



\begin{thebibliography}{99}



\RBibitem{balk:c1}

\paper К проблеме уравнений смешанного типа

\by Лаврентьев~М.~А., Бицадзе~А.~В.

\jour Докл. АН СССР

\yr 1950

\vol 70

\issue 3

\pages 373--376

\mathscinet{http://www.ams.org/mathscinet-getitem?mr=35371}



\RBibitem{balk:c2}

\by Бицадзе~А.~В.

\book Уравнения смешанного типа

\publaddr М.

\publ АН СССР

\yr 1959

\totalpages 164



\RBibitem{balk:c3}

\by Чибрикова~Л.~И.

\paper Новый метод решения одной краевой задачи смешанного типа

\serial Учен. зап. Казанск. ун-та

\yr 1957

\vol 117

\issue 9

\pages 44--47

\publ Изд-во Казанск. ун-та

\publaddr Казань



\RBibitem{balk:c4}

\by Чибрикова~Л.~И.

\paper К решению краевой задачи Трикоми для уравнения

${u_{xx}+\operatorname{sgn}yu_{yy}=0}$

\serial Учен. зап. Казанск. ун-та

\yr 1957

\vol 117

\issue 9

\pages 48--51

\publ Изд-во Казанск. ун-та

\publaddr Казань



\RBibitem{balk:c5}

\by Лернер~М.~Е., Пулькин~С.~П.

\paper О~единственности решений задач с условиями Франкля и Трикоми для общего уравнения Лаврентьева--Бицадзе

\jour Дифференц. уравнения

\yr 1966

\vol 2

\issue 9

\pages 1255--1263

\mathnet{http://mi.mathnet.ru/de576}

\mathscinet{http://www.ams.org/mathscinet-getitem?mr=206518}

\zmath{https://zbmath.org/?q=an:0152.09803}





\RBibitem{balk:c6}

\by Крикунов~Ю.~М.

\paper К~задаче Трикоми для уравнения Лаврентьева--Бицадзе

\jour Изв. вузов. Матем.

\yr 1974

\issue 2

\pages 76--81

\mathnet{http://mi.mathnet.ru/ivm6523}

\mathscinet{http://www.ams.org/mathscinet-getitem?mr=352720}

\zmath{https://zbmath.org/?q=an:0282.35062}





\RBibitem{balk:c7}

\by Моисеев~Е.~И.

\paper О~задаче Трикоми для уравнения Лаврентьева--Бицадзе

\jour Дифференц. уравнения

\yr 1981

\vol 17

\issue 2

\pages 325--338

\mathnet{http://mi.mathnet.ru/de4197}

\mathscinet{http://www.ams.org/mathscinet-getitem?mr=606823}

\zmath{https://zbmath.org/?q=an:0456.35059}



\RBibitem{balk:c8}

\by Казиев~В.~М.

\paper Задача Трикоми для нагруженного уравнения Лаврентьева--Бицадзе

\jour Дифференц. уравнения

\yr 1979

\vol 15

\issue 1

\pages 173--175

\mathnet{http://mi.mathnet.ru/de3610}

\mathscinet{http://www.ams.org/mathscinet-getitem?mr=524950}

\zmath{https://zbmath.org/?q=an:0398.35068}



\RBibitem{balk:c9}

\by Сабитов~К.~Б.

\paper О~задаче Трикоми для уравнения Лаврентьева--Бицадзе со спектральным параметром

\jour Дифференц. уравнения

\yr 1986

\vol 22

\issue 11

\pages 1977--1984

\mathnet{http://mi.mathnet.ru/de6041}

\mathscinet{http://www.ams.org/mathscinet-getitem?mr=876371}



\RBibitem{balk:c10}

\by Солдатов~А.~П.

\paper Задачи типа Дирихле для уравнения Лаврентьева--Бицадзе

\jour Дифференц. уравнения

\yr 1994

\vol 30

\issue 11

\pages 2001--2009

\mathnet{http://mi.mathnet.ru/de8498}

\mathscinet{http://www.ams.org/mathscinet-getitem?mr=1348291}

%\transl

%\jour Differ. Equ.

%\yr 1994

%\vol 30

%\issue 11

%\pages 1846--1853



\RBibitem{balk:c11}

\by Репин~О.~А.

\paper Задача Трикоми для уравнения смешанного типа в~области, эллиптическая часть которого~-- полуполоса

\jour Дифференц. уравнения

\yr 1996

\vol 32

\issue 4

\pages 565--567

\mathnet{http://mi.mathnet.ru/de8986}

\mathscinet{http://www.ams.org/mathscinet-getitem?mr=1436997}

%\transl

%\jour Differ. Equ.

%\yr 1996

%\vol 32

%\issue 4

%\pages 571--574



\RBibitem{balk:c12}

\by Лернер~М.~Е., Репин~О.~А.

\paper Об одной задаче с двумя нелокальными краевыми условиями для уравнения смешанного типа

\jour Сиб. матем. журн.

\yr 1999

\vol 40

\issue 6

\pages 1260--1275

\mathnet{http://mi.mathnet.ru/smj64}

\mathscinet{http://www.ams.org/mathscinet-getitem?mr=1741080}

\zmath{https://zbmath.org/?q=an:1022.35030}

%\transl

%\jour Siberian Math. J.

%\yr 1999

%\vol 40

%\issue 6

%\pages 1064--1078

%\crossref{10.1007/BF02677530}

%\isi{http://gateway.isiknowledge.com/gateway/Gateway.cgi?GWVersion=2&SrcApp=PARTNER_APP&SrcAuth=LinksAMR&DestLinkType=FullRecord&DestApp=ALL_WOS&KeyUT=000084386400006}



\RBibitem{balk:c13}

\by Нахушева~З.~А.

\paper О~задаче Трикоми для уравнения Лаврентьева--Бицадзе с~нелокальным условием сопряжения

\jour Дифференц. уравнения

\yr 2005

\vol 41

\issue 10

\pages 1426--1428

\mathnet{http://mi.mathnet.ru/de11382}

\mathscinet{http://www.ams.org/mathscinet-getitem?mr=2241472}

%\transl

%\jour Differ. Equ.

%\yr 2005

%\vol 41

%\issue 10

%\pages 1505--1508

%\crossref{10.1007/s10625-005-0307-y}



\RBibitem{balk:c14}

\by Ладыженская~О.~А.

\paper Об одном способе приближённого решения задачи Лаврентьева--Бицадзе

\jour УМН

\yr 1954

\vol 9

\issue 4(62)

\pages 187--189

\mathnet{http://mi.mathnet.ru/umn8126}

\mathscinet{http://www.ams.org/mathscinet-getitem?mr=70033}

\zmath{https://zbmath.org/?q=an:0056.31905}



\RBibitem{balk:c15}

\by Халилов~З.~И.

\paper Решение задачи для уравнения смешанного типа методом сеток

\jour Докл. АН АзССР

\yr 1953

\vol 9

\issue 4

\pages 189--190



\RBibitem{balk:c16}

\by Карманов~В.~Г.

\paper Об одной граничной задаче для уравнения смешанного типа

\jour Докл. АН СССР

\yr 1954

\vol 95

\issue 3

\pages 439--442

\mathscinet{http://www.ams.org/mathscinet-getitem?mr=64994}

\zmath{https://zbmath.org/?q=an:0055.32901}



\RBibitem{balk:c17}

\by Филиппов~А.~Ф.

\paper О разностном методе решения задачи Трикоми

\jour Изв. АН СССР. Сер. матем.

\yr 1957

\vol 21

\issue 1

\pages 73--88

\mathnet{http://mi.mathnet.ru/izv4005}

\mathscinet{http://www.ams.org/mathscinet-getitem?mr=90133}

\zmath{https://zbmath.org/?q=an:0079.12001}



\Bibitem{balk:c18}

\paper On Difference Methods for the Solution of a Tricomi Problem 

\by Ogawa~H.

\jour Trans. Amer. Math. Soc.

\vol 100

\yr 1961

\issue 3

\pages 404--424

\crossref{10.2307/1993522}



\RBibitem{balk:c19}

\by Сабитов~К.~Б., Мугафаров~М.~Ф.

\paper Экстремальные свойства решений разностной задачи Трикоми для одной сеточной системы уравнений смешанного типа и их применения

\jour Изв. вузов. Матем.

\yr 2005

\issue 4

\pages 56--69

\mathnet{http://mi.mathnet.ru/ivm38}

\mathscinet{http://www.ams.org/mathscinet-getitem?mr=2180685}

\zmath{https://zbmath.org/?q=an:1118.35015}

\elib{http://elibrary.ru/item.asp?id=9083825}

%\transl

%\jour Russian Math. (Iz. VUZ)

%\yr 2005

%\vol 49

%\issue 4

%\pages 53--66



\RBibitem{balk:c20}

\paper Два газодинамических приложения краевой задачи Лаврентьева--Бицадзе

\by Франкль~Ф.~И.

\jour Вестник ЛГУ. Сер. матем., мех. и астр.

\vol 6

\yr 1951

\issue 11

\pages 3--7





\RBibitem{balk:c21}

\by Самарский~A.~A.

\book Теория разностных схем

\yr 1983

\publ Наука

\publaddr М.

\mathscinet{http://www.ams.org/mathscinet-getitem?mr=721476}

\zmath{https://zbmath.org/?q=an:0381.76051}

\totalpages 616



\end{thebibliography}





\newpage



\makeentitle





\AdditionalContent{Competing interests}{We have no competing interests.} 

\AdditionalContent{Authors' contributions and responsibilities}{Each author has participated in the article concept development and in the manuscript writing.

The authors are absolutely responsible for submitting the final manuscript in print.

Each author has approved the final version of manuscript.} 

\AdditionalContent{Funding}{The research has not had any sponsorship.}



\newpage



\begin{thebibliography}{99}



\Bibitem{balk:c1:en}

\lang In Russian

\paper On the problem of equations of mixed type

\by Lavrent'ev~M.~A., Bitsadze~A.~V.

\jour Doklady Akad. Nauk SSSR

\yr 1950

\vol 70

\issue 3

\pages 373--376

\mathscinet{http://www.ams.org/mathscinet-getitem?mr=35371}



\Bibitem{balk:c2:en}

\lang In Russian

\by Bitsadze~A.~V. 

\book Uravneniia smeshannogo tipa \rm [Equations of mixed type]

\publaddr Moscow

\publ USSR Academy of Sciences Publ.

\yr 1959 

\totalpages 164





\Bibitem{balk:c3:en}

\lang In Russian

\by Chibrikova~L.~I.

\paper New method of solving one boundary value problem of the mixed type

\yr 1957

\vol 117

\issue 9

\pages 44--47

\serial Uchenye Zapiski Kazanskogo Universiteta

\publ Kazan State University

\publaddr Kazan







\Bibitem{balk:c4:en}

\lang In Russian

\by Chibrikova~L.~I.

\paper To the solution of the Tricomi boundary value problem for equation

${u_{xx}+\operatorname{sgn}yu_{yy}=0}$

\yr 1957

\vol 117

\issue 9

\pages 48--51

\serial Uchenye Zapiski Kazanskogo Universiteta

\publ Kazan State University

\publaddr Kazan



\Bibitem{balk:c5:en}

\lang In Russian

\by Lerner~M.~E., Pul'kin~S.~P.

\paper Uniqueness of a solution of problems with Frankl' and Tricomi conditions for the general Lavrent'ev--Bitsadze equation

\jour Differ. Uravn.

\yr 1966

\vol 2

\issue 9

\pages 1255--1263





\Bibitem{balk:c6:en}

\lang In Russian

\by Krikunov~Yu.~M.

\paper On the Tricomi problem for the Lavrent'ev--Bitsadze equation

\jour Izv. Vyssh. Uchebn. Zaved. Mat.

\yr 1974

\issue 2

\pages 76--81





\Bibitem{balk:c7:en}

\lang In Russian

\by E.~I.~Moiseev

\paper The Tricomi problem for the Lavrent'ev--Bitsadze equation

\jour Differ. Uravn.

\yr 1981

\vol 17

\issue 2

\pages 325--338



\Bibitem{balk:c8:en}

\lang In Russian

\by Kaziev~V.~M.

\paper The Tricomi problem for a loaded Lavrent'ev--Bitsadze equation

\jour Differ. Uravn.

\yr 1979

\vol 15

\issue 1

\pages 173--175



\Bibitem{balk:c9:en}

\lang In Russian

\by Sabitov~K.~B.

\paper The Tricomi problem for the Lavrent'ev--Bitsadze equation with a spectral parameter

\jour Differ. Uravn.

\yr 1986

\vol 22

\issue 11

\pages 1977--1984





\Bibitem{balk:c10:en}

\by Soldatov~A.~P.

\paper Dirichlet-type problems for the Lavrent'ev--Bitsadze equation

\jour Differ. Equ.

\yr 1994

\vol 30

\issue 11

\pages 1846--1853

\mathscinet{http://www.ams.org/mathscinet-getitem?mr=1348291}



\Bibitem{balk:c11:en}

\by Repin~O.~A.

\paper The Tricomi problem for an equation of mixed type in a domain whose elliptic part is a half-strip

\mathscinet{http://www.ams.org/mathscinet-getitem?mr=1436997}

\jour Differ. Equ.

\yr 1996

\vol 32

\issue 4

\pages 571--574



\Bibitem{balk:c12:en}

\paper A problem with two nonlocal boundary conditions for a mixed-type equation

\by Lerner~M.~E.,  Repin~O.~A.

\jour Siberian Math. J.

\yr 1999

\vol 40

\issue 6

\pages 1064--1078

\crossref{10.1007/BF02677530}

\isi{http://gateway.isiknowledge.com/gateway/Gateway.cgi?GWVersion=2&SrcApp=PARTNER_APP&SrcAuth=LinksAMR&DestLinkType=FullRecord&DestApp=ALL_WOS&KeyUT=000084386400006}



\Bibitem{balk:c13:en}

\jour Differ. Equ.

\yr 2005

\vol 41

\issue 10

\pages 1505--1508

\crossref{10.1007/s10625-005-0307-y}

\paper On the Tricomi Problem for the Lavrent'ev--Bitsadze Equation with a Nonlocal Transmission Condition

\by  Nakhusheva~Z.~A.



\Bibitem{balk:c14:en}

\lang In Russian

\by Ladyzhenskaya~O.~A.

\paper On a~method of approximate solution of the Lavrent'ev--Bicadze problem

\jour Uspekhi Mat. Nauk

\yr 1954

\vol 9

\issue 4(62)

\pages 187--189



\Bibitem{balk:c15:en}

\lang In Russian

\by Khalilov~Z.~I.

\paper Solution of the problem for a mixed-type equation by the grid method

\jour Doklady Akad. Nauk AzSSR

\yr 1953

\vol 9

\issue 4

\pages 189--190





\Bibitem{balk:c16:en}

\lang In Russian

\by Karmanov~V.~G.

\paper On a boundary problem for an equation of mixed type

\jour Doklady Akad. Nauk SSSR 

\yr 1954

\vol 95

\issue 3

\pages 439--442





\Bibitem{balk:c17:en}

\lang In Russian

\by Filippov~A.~F.

\paper On the application of the method of finite differences to the solution of the problem of Tricomi

\jour Izv. Akad. Nauk SSSR Ser. Mat.

\yr 1957

\vol 21

\issue 1

\pages 73--88





\Bibitem{balk:c18:en}

\paper On Difference Methods for the Solution of a Tricomi Problem 

\by Ogawa~H.

\jour Trans. Amer. Math. Soc.

\vol 100

\yr 1961

\issue 3

\pages 404--424

\crossref{10.2307/1993522}



\Bibitem{balk:c19:en}

\by Sabitov~K.~B., Mugafarov~M.~F.

\paper Extremal properties of solutions of the Tricomi difference problem for a difference system of equations of mixed type and their applications

\mathscinet{http://www.ams.org/mathscinet-getitem?mr=2180685}

\zmath{https://zbmath.org/?q=an:1118.35015}

\jour Russian Math. (Iz. VUZ)

\yr 2005

\vol 49

\issue 4

\pages 53--66



\Bibitem{balk:c20:en}

\lang In Russian

\paper Two gas-dynamical applications of the Lavrentiev–Bitsadze boundary value problem

\by Frankl~F.~I.

\jour Vestnik Leningrad. Univ. Ser. Mat. Meh. Astronom

\vol 6

\yr 1951

\issue 11

\pages 3--7







\Bibitem{balk:c21:en}

\lang In Russian

\by Samarskii~A.~A.

\book Teoriia raznostnykh skhem \rm [Theory of difference schemes] 

\yr 1983

\publ Nauka

\publaddr Moscow

\totalpages 616



\end{thebibliography}



\renewcommand{\baselinestretch}{0.9}



%\end{document}

